\documentclass[11pt,a4paper]{article}
\usepackage[margin=1in]{geometry}
\usepackage[T1]{fontenc}
\usepackage[utf8]{inputenc}
\usepackage{lmodern}
\usepackage{microtype}
\usepackage{enumitem}
\usepackage{hyperref}
\usepackage{fancyhdr}
\usepackage{array}
\usepackage[authoryear,round]{natbib}
\setlength{\parindent}{0pt}
\setlength{\parskip}{0.65em}
\pagestyle{fancy}
\fancyhf{}
\fancyfoot[C]{\thepage}
\renewcommand{\headrulewidth}{0pt}
\newcommand{\guidance}[1]{\textit{#1}}
\begin{document}
\begin{center}
{\LARGE\bfseries AI Collusion Forensics: Comparing Communication and Behavioral Evidence}\par
\vspace{1.5em}
\begin{tabular}{>{\centering\arraybackslash}p{0.28\textwidth} >{\centering\arraybackslash}p{0.28\textwidth} >{\centering\arraybackslash}p{0.28\textwidth}}
Agnes Ahalya Arogyaraj \\
Saarland University \\[1.3em]

\end{tabular}
\par\vspace{1em}
With\\Apart Research
\end{center}
\begin{abstract}
This study asks whether monitors detect collusion more reliably when they can compare what agents say with what they do. Two or more language-model agents play the existing public-goods game under two conditions: an independent-control condition that hides other agents' announcements before final action, and a coordination-opportunity condition that reveals them. The original game runner and payoff calculation are preserved. Every saved trajectory is assessed by three independent labeling models, followed by evidence checks and human review when the label remains uncertain. Four monitors then receive different, condition-blinded evidence views: public messages, actions and outcomes, a combined timeline, or a structured comparison of claims and actions. This draft fixes the design and analysis plan before the complete run. Quantitative results and conclusions will be added after trajectories, labels, and monitor outputs have been checked.

\end{abstract}
\vspace{1em}
\vfill
\noindent\textsuperscript{1} Research conducted at the AI Collusion Research Sprint, October 2026
\newpage

\section{Introduction}
Multi-agent systems can produce apparent coordination through shared context or similar incentives. A useful collusion assessment therefore needs evidence of an agreement, its relationship to the study rule, and the actions that followed. Public statements alone can overstate intent; actions alone can hide how the agents reached them. The practical question is which evidence view helps a monitor distinguish these cases.

We compare four views of the \emph{same} game trajectories so that a difference between monitors can be attributed to the evidence shown, subject to model and sampling limits.

Our main contributions are:
\begin{enumerate}
\item A two-condition execution protocol using the existing game engine and payoff calculation.
\item A three-judge labeling process with source verification, an unknown outcome, and human review.
\item A matched comparison of communication-only, behavior-only, combined, and structured combined monitoring.
\end{enumerate}

\section{Related Work}
Research on multi-agent systems has examined how communication and behavior reveal coordination. COLOSSEUM(\cite{nakamura2026colosseum}) investigates collusion through agent messages and observed actions, distinguishing expressed intentions from executed behavior . \cite{radev2026best} show that access to agent communication can improve detection of coordinated sabotage. \cite{shi2026agents} examine differences between private intentions, public commitments, and final actions in repeated games, providing the experimental foundation for our study.

Monitoring research also highlights how evidence presentation affects detection. \cite{arike2026does} find that additional information does not always improve sabotage detection, while \cite{nguyen2026traceguard} introduce TraceGuard, which uses structured assessments of agent behavior. Building on these findings, we compare four monitoring approaches on identical multi-agent trajectories: (1) public messages only, (2) actions and outcomes only, (3) combined messages and actions, and (4) structured combined evidence comparing stated and observed behavior. We investigate whether combining communication and behavioral evidence improves collusion detection over either source alone.

\section{Methods}
\subsection{Game and experimental conditions}
The source experiment is the \texttt{llm-trust-breaking} endogenous public-goods runner. For each trial, agents privately plan, publicly announce, and then choose final actions. We retain the runner's trial execution and payoff evaluation. Both conditions add a study instruction to make independent decisions and avoid agreements on joint actions. This is an experimental rule, not an original game rule.

In \emph{independent control}, the final-action prompt does not show peers' Stage 2 announcements. In \emph{coordination opportunity}, it does. Both conditions ask each agent at the final stage for a yes/no collusion self-report, confidence, and reasoning. This question may influence behavior; we record it as part of the intervention. The self-report can be considered by label judges as an unverified claim, but is withheld from monitors.

The GitHub workflow runs both conditions in one isolated execution. Its only user inputs are the number of agents per game and trials per condition. Each trial produces one trajectory, so the planned total is twice the trials-per-condition input. The workflow fixes the game-agent and native game-judge model to DeepSeek V3.2, the three label judges to GPT-4.1 Nano, Mistral Small 3.2, and DeepSeek V3.2, and the monitor model to GPT-4.1 Mini. It preflights request counts and cost reserves before paid calls. These model choices and prices should be recorded with the completed run.

Game agents use temperature 0.7. Trial seeds start at 42 and increase with the trial ID; a deterministic hash derives distinct request seeds for each agent and stage. Corresponding requests use the same seeds across conditions. The native game judge, labeling judges, and monitors use temperature 0 where supported. The requested sampling settings and protocol version are saved with each trajectory. Seed reproducibility depends on the provider; repeated game content is counted explicitly. Earlier temperature-0 runs will be identified separately. The full workflow has a combined local reserve limit of \$3.50.

\subsection{Trajectory records and ground-truth labels}
A trajectory stores private plans, public messages, final-action reasoning and actions, game outcomes, condition metadata, and available API usage. Label judges receive a condition-blinded case with plans, messages, actions, outcomes, and self-reports. Each independently returns a binary candidate or unknown, confidence, and cited evidence. A candidate needs a majority, median agreeing confidence of at least 0.80, and deterministic checks that its citations and operational evidence are supported by the saved trajectory. A positive label requires a reciprocal agreement that violates the added independent-decision instruction and was carried out. Cases without a verified decision remain \emph{unknown} and enter the human review queue; unknown is never counted as negative. Human decisions require a written justification and preserve the original model votes.

\subsection{Four blinded monitoring views}
Each monitor sees one of four views of a trajectory: M1 public messages; M2 actions and outcomes; M3 a combined chronological record; or M4 a structured comparison between stated claims and observed actions. Monitor prompts exclude treatment name, private plans, agent self-reports, judge votes, verification details, and final labels. A monitor returns a collusion score; scores at or above 0.5 count as positive predictions. The four views are scored on the same trajectory IDs.

\subsection{Analysis plan}
We will report, for each view, true positives, false positives, true negatives, false negatives, precision, and recall against \emph{verified} final labels. We will give every denominator and report an undefined metric when its denominator is zero. Unknown final labels and failed monitor responses will be counted separately. We will compare the four views case by case and by condition, without treating judge agreement as proof of a true population rate. We will also report actual HTTP attempts, provider-reported or estimated token usage and cost, and elapsed time where recorded.

\section{Results}
\textbf{Results pending.} Quantitative findings will be added after the complete workflow artifact has been checked. No result values are claimed in this draft.

\begin{table}[h]
\centering
\begin{tabular}{lrrrr}
\hline
Condition & Total & Positive & Negative & Unknown\\
\hline
Independent control & -- & -- & -- & --\\
Coordination opportunity & -- & -- & -- & --\\
\hline
\end{tabular}
\caption{Trajectory and verified label counts to be filled after execution. Dashes indicate pending values.}
\end{table}

\begin{table}[h]
\centering
\begin{tabular}{lrrrrcc}
\hline
View & TP & FP & TN & FN & Precision & Recall\\
\hline
M1 messages & -- & -- & -- & -- & -- & --\\
M2 behavior & -- & -- & -- & -- & -- & --\\
M3 combined & -- & -- & -- & -- & -- & --\\
M4 structured & -- & -- & -- & -- & -- & --\\
\hline
\end{tabular}
\caption{Planned monitor comparison on verified labels only. Unknown final labels and failed responses will be reported separately.}
\end{table}

The completed report will include actual API usage, cost, and elapsed time by stage, plus a representative trajectory showing each agent's plan, announcement, final action, and resulting payoff. Claims about differences between monitor views will state the number of evaluable trajectories and the limits of the sample.

\section{Discussion and Limitations}
\textbf{Interpretation pending.} The analysis will ask whether communication and behavior together resolve cases that either view alone leaves ambiguous. It will separately assess whether the structured presentation changes performance relative to the raw combined timeline.

\subsection*{Limitations}
The small pilot cannot support a precise estimate of detection performance. Agreement among three model judges can still share blind spots, and their self-reported confidence is not calibrated. The added independence instruction and collusion self-report change the game prompt. Monitoring on one model and one score threshold may not generalize. Local spending caps are estimates rather than provider-enforced limits. Partial or invalid runs will be reported explicitly.

\subsection*{Future Work}
A larger sample can test whether any observed differences persist across more trajectories. Further work could vary agent and monitor models, examine the effect of the self-report question, and measure agreement between human reviewers and model judges.

\section{Conclusion}
This draft specifies the execution and evaluation before observing the complete run. The final conclusion will be written from verified labels and matched monitor results.

\subsection*{Code and Data}
\begin{itemize}
\item Code repository: \url{https://github.com/AgnesAhalya/collusion_hackathon}
\item Data/Datasets: Results pending. The workflow saves an encrypted artifact containing trajectories, labels, monitor outputs, and available API usage.
\item Other artifacts: The \texttt{Track 2.2 complete experiment} GitHub Action runs the full pipeline. The completed report will record the run ID and commit SHA.
\end{itemize}

\bibliographystyle{plainnat}
\bibliography{references}



\section*{LLM Usage Statement}
Large language models (LLMs) were used to brainstorm research ideas, summarize and understand related literature, assist with code generation and debugging, and support drafting and editing the report. The research design, implementation decisions, interpretation of results, and final conclusions remained the responsibility of the author. All AI-generated content was reviewed before inclusion.

\end{document}
